% !TeX program = xelatex
\documentclass[11pt,a4paper]{article}

% -------------------------------------------------
% Packages
% -------------------------------------------------
\usepackage{kotex}
\usepackage{fontspec}
\usepackage{amsmath,amssymb,bm,mathtools}
\usepackage[left=24mm,right=24mm,top=23mm,bottom=25mm]{geometry}
\usepackage{booktabs,array}
\usepackage{enumitem}
\usepackage{setspace}
\usepackage{xcolor}
\usepackage{colortbl}
\usepackage{tikz}
\usepackage{pifont}
\usetikzlibrary{positioning,arrows.meta,calc,fit,decorations.pathreplacing}
\usepackage[most]{tcolorbox}
\usepackage{listings}
\usepackage{hyperref}
\usepackage{fancyhdr}

\setstretch{1.14}
\setlength{\parindent}{0pt}
\setlength{\parskip}{0.55em}

% -------------------------------------------------
% Colors and boxes
% -------------------------------------------------
\definecolor{deepblue}{RGB}{34,71,120}
\definecolor{softblue}{RGB}{239,246,253}
\definecolor{softgray}{RGB}{246,247,249}
\definecolor{softgreen}{RGB}{239,249,242}
\definecolor{softyellow}{RGB}{255,249,229}
\definecolor{codegray}{RGB}{248,248,248}

\newtcolorbox{keybox}[1][]{
  colback=softblue,
  colframe=deepblue,
  boxrule=0.7pt,
  arc=2mm,
  title=#1,
  fonttitle=\bfseries
}

\newtcolorbox{thinkbox}[1][]{
  colback=softyellow,
  colframe=orange!70!black,
  boxrule=0.6pt,
  arc=2mm,
  title=#1,
  fonttitle=\bfseries
}

\newtcolorbox{remarkbox}[1][]{
  colback=softgray,
  colframe=black!45,
  boxrule=0.5pt,
  arc=2mm,
  title=#1,
  fonttitle=\bfseries
}

\newtcolorbox{mlbox}[1][]{
  colback=softgreen,
  colframe=green!45!black,
  boxrule=0.6pt,
  arc=2mm,
  title=#1,
  fonttitle=\bfseries
}

% -------------------------------------------------
% Code style
% -------------------------------------------------
\lstdefinestyle{pythonstyle}{
  language=Python,
  backgroundcolor=\color{codegray},
  basicstyle=\ttfamily\small,
  keywordstyle=\bfseries,
  commentstyle=\itshape\color{black!60},
  stringstyle=\color{black!75},
  numbers=left,
  numberstyle=\tiny\color{black!55},
  numbersep=8pt,
  frame=single,
  rulecolor=\color{black!25},
  breaklines=true,
  showstringspaces=false,
  tabsize=4,
  columns=fullflexible,
  keepspaces=true
}
\lstset{style=pythonstyle}

% -------------------------------------------------
% Math commands
% -------------------------------------------------
\newcommand{\R}{\mathbb{R}}
\newcommand{\grad}{\nabla}
\newcommand{\vect}[1]{\bm{#1}}
\newcommand{\T}{^{\mathsf T}}
\newcommand{\loss}{\mathcal{L}}
\newcommand{\sigmoid}{\sigma}
\newcommand{\pd}[2]{\frac{\partial #1}{\partial #2}}
\newcommand{\softmax}{\operatorname{softmax}}

% -------------------------------------------------
% Header / footer
% -------------------------------------------------
\pagestyle{fancy}
\fancyhf{}
\lhead{Introduction to Deep Learning and Practice}
\rhead{Multi-class Classification: Answers}
\cfoot{\thepage}
\renewcommand{\headrulewidth}{0.4pt}
\setlength{\headheight}{14pt}

\hypersetup{
  colorlinks=true,
  linkcolor=deepblue,
  urlcolor=deepblue
}

% =================================================

% source tag of a problem
\newcommand{\src}[1]{{\unskip\nobreak\hfil\penalty50\hskip1em\hbox{}\nobreak\hfill\mbox{\small\color{black!50}[#1]}}}
% a problem block
\newlist{problems}{enumerate}{1}
\setlist[problems]{label=\textbf{\arabic*.},leftmargin=2.2em,itemsep=0.2em}
% =================================================
% a solution block
\newtcolorbox{sol}{breakable,enhanced,colback=softgray,colframe=black!35,boxrule=0.4pt,arc=1mm,
  left=2mm,right=2mm,top=1mm,bottom=1mm,before upper={\textbf{풀이.}\ }}
\setlist[problems]{label=\textbf{\arabic*.},leftmargin=2.2em,itemsep=0.9em}
\begin{document}

\begin{center}
{\LARGE\bfseries 연습문제 정답과 풀이}\\[0.35em]
{\large 다중분류 (MNIST) $\cdot$ 딥러닝개론및연습 2026년 2학기}
\end{center}

강의노트 끝의 연습문제 8개의 정답과 풀이이다. 문제를 다시 적고 그 아래에 풀이를 쓴다.

\begin{problems}

% ---------------------------------------------------------------- 1
\item MNIST network에서 $l=784$, $m=100$, $n=10$, 미니배치 크기 $h=100$이다.
$X$, $W^{(1)}$, $B^{(1)}$, $A$, $Z$, $W^{(2)}$, $B^{(2)}$, $U$, $Y$, $T$의 shape을 쓰고,
학습해야 하는 parameter의 총 개수를 구하라.

\begin{sol}
미니배치의 데이터 하나가 한 행이므로 행의 수는 $h$이고, 열의 수는 그 층의 node 수이다.
\begin{center}
\begin{tabular}{ll@{\qquad}ll}
\toprule
변수 & shape & 변수 & shape \\
\midrule
$X$ (입력) & $h\times l=100\times784$ & $W^{(2)}$ & $m\times n=100\times10$ \\
$W^{(1)}$ & $l\times m=784\times100$ & $B^{(2)}$ & $1\times n=1\times10$ \\
$B^{(1)}$ & $1\times m=1\times100$ & $U=ZW^{(2)}+B^{(2)}$ & $h\times n=100\times10$ \\
$A=XW^{(1)}+B^{(1)}$ & $h\times m=100\times100$ & $Y=\softmax(U)$ & $h\times n=100\times10$ \\
$Z=h(A)$ & $h\times m=100\times100$ & $T$ (one-hot) & $h\times n=100\times10$ \\
\bottomrule
\end{tabular}
\end{center}
$B^{(1)}$, $B^{(2)}$는 $1\times m$, $1\times n$ 행벡터이고, 더할 때 $h$개의 행에 broadcast된다.
\begin{center}
\end{center}
학습하는 parameter는 $W^{(1)},B^{(1)},W^{(2)},B^{(2)}$이다.
\[
784\cdot100+100+100\cdot10+10=78400+100+1000+10=\boxed{79510}.
\]
($X$, $A$, $Z$, $U$, $Y$, $T$는 데이터에 따라 정해지는 값이고 parameter가 아니다.
또 parameter의 개수는 미니배치 크기 $h$와 관계없다.)
\end{sol}

% ---------------------------------------------------------------- 2
\item $U=(1,2,3)$일 때 $Y=\softmax(U)$를 소수 셋째 자리까지 구하라.
또 임의의 상수 $c$에 대하여 $\softmax(U+c)=\softmax(U)$임을 보이라.

\begin{sol}
$e^1\approx2.718$, $e^2\approx7.389$, $e^3\approx20.086$이고 합은 $30.193$이다. 따라서
\[
Y\approx\Bigl(\frac{2.718}{30.193},\ \frac{7.389}{30.193},\ \frac{20.086}{30.193}\Bigr)
=\boxed{(0.090,\ 0.245,\ 0.665)}.
\]
세 값의 합은 $1$이다. 점수가 $1$씩 커질 때마다 확률은 약 $e\approx2.7$배씩 커진다.

$U+c=(u_1+c,\ldots,u_n+c)$이면 $e^{u_i+c}=e^c e^{u_i}$이므로
\[
\softmax(U+c)_i=\frac{e^{c}e^{u_i}}{\sum_{j}e^{c}e^{u_j}}
=\frac{e^{c}e^{u_i}}{e^{c}\sum_{j}e^{u_j}}
=\frac{e^{u_i}}{\sum_{j}e^{u_j}}=\softmax(U)_i .
\]
분자와 분모의 $e^c$가 약분된다.
코드의 \texttt{U = U - np.max(U, axis=1, keepdims=True)}는 $c=-\max_j u_j$로 둔 것이다.
값은 그대로이면서 지수의 최댓값이 $e^0=1$이 되어 overflow가 생기지 않는다.
\end{sol}

% ---------------------------------------------------------------- 3
\item 정답이 $t=(0,0,1)$이고 network의 출력이 $Y=(0.2,\,0.3,\,0.5)$일 때
cross entropy 손실 $l$과 $\pd{l}{U}$를 구하라.
경사하강법은 세 점수 $u_1,u_2,u_3$를 각각 어느 방향으로 움직이는가?

\begin{sol}
$t$에서 $1$인 성분은 셋째뿐이므로
\[
l=-\sum_i t_i\log y_i=-\log y_3=-\log0.5=\log2\approx\boxed{0.693}.
\]
softmax와 cross entropy를 묶은 계층의 역전파는 $\pd{l}{U}=Y-t$이므로
\[
\pd{l}{U}=(0.2-0,\ 0.3-0,\ 0.5-1)=\boxed{(0.2,\ 0.3,\ -0.5)}.
\]
경사하강법은 gradient의 반대 방향으로 움직인다 ($u_i\leftarrow u_i-\alpha\,\pd{l}{u_i}$).
\begin{itemize}[leftmargin=1.5em,itemsep=0pt]
\item $u_1$, $u_2$: gradient가 양수이므로 \textbf{줄어든다}. 틀린 class의 점수를 낮춘다.
\item $u_3$: gradient가 음수이므로 \textbf{커진다}. 정답 class의 점수를 높인다.
\end{itemize}
움직이는 양은 $y_i-t_i$에 비례한다. 즉 정답과 많이 다를수록 많이 움직인다.
또 세 성분의 합이 $0.2+0.3-0.5=0$이므로, 모든 점수에 같은 값을 더하는 방향으로는 움직이지 않는다 (문제 2와 같은 이유).
\end{sol}

% ---------------------------------------------------------------- 4
\item Affine 계층 $A=XW+B$에서
\[
X=(1,\ 2),\qquad
W=\begin{pmatrix}1&0&-1\\2&1&0\end{pmatrix},\qquad
B=(0,\ 1,\ 0),\qquad
\pd{l}{A}=(1,\ -1,\ 2)
\]
일 때 $A$, $\pd{l}{X}$, $\pd{l}{W}$, $\pd{l}{B}$를 구하라.

\begin{sol}
순전파:
\[
XW=(1\cdot1+2\cdot2,\ 1\cdot0+2\cdot1,\ 1\cdot(-1)+2\cdot0)=(5,\ 2,\ -1),\qquad
A=XW+B=\boxed{(5,\ 3,\ -1)}.
\]
역전파 ($\pd{l}{X}=\pd{l}{A}W\T$, $\pd{l}{W}=X\T\pd{l}{A}$, $\pd{l}{B}=\pd{l}{A}$):
\[
\pd{l}{X}=(1,\ -1,\ 2)\begin{pmatrix}1&2\\0&1\\-1&0\end{pmatrix}
=(1-2,\ 2-1)=\boxed{(-1,\ 1)},
\]
\[
\pd{l}{W}=\begin{pmatrix}1\\2\end{pmatrix}(1,\ -1,\ 2)
=\boxed{\begin{pmatrix}1&-1&2\\2&-2&4\end{pmatrix}},\qquad
\pd{l}{B}=\boxed{(1,\ -1,\ 2)}.
\]
shape 확인: $\pd{l}{X}$는 $X$와 같은 $1\times2$, $\pd{l}{W}$는 $W$와 같은 $2\times3$, $\pd{l}{B}$는 $B$와 같은 $1\times3$이다.
\end{sol}

% ---------------------------------------------------------------- 5
\item softmax의 Jacobian $J=\left(\pd{y_j}{u_i}\right)$가
\[
J=\operatorname{diag}(y_1,\ldots,y_n)-Y\T Y
\]
로 쓸 수 있음을 보이라. ($Y$는 $1\times n$ 행벡터)

\begin{sol}
$S=\sum_k e^{u_k}$라 하면 $y_j=e^{u_j}/S$이고 $\pd{S}{u_i}=e^{u_i}$이다.
몫의 미분법으로
\[
\pd{y_j}{u_i}
=\frac{\delta_{ij}e^{u_j}\,S-e^{u_j}e^{u_i}}{S^2}
=\delta_{ij}\frac{e^{u_j}}{S}-\frac{e^{u_j}}{S}\cdot\frac{e^{u_i}}{S}
=\delta_{ij}\,y_j-y_iy_j .
\]
($\delta_{ij}$는 $i=j$이면 $1$, 아니면 $0$.)
첫째 항을 $(i,j)$ 성분으로 갖는 행렬은 $\operatorname{diag}(y_1,\ldots,y_n)$이다.
둘째 항 $y_iy_j$는 열벡터 $Y\T$ ($n\times1$)와 행벡터 $Y$ ($1\times n$)의 곱 $Y\T Y$ ($n\times n$)의 $(i,j)$ 성분이다. 따라서
\[
J=\operatorname{diag}(y_1,\ldots,y_n)-Y\T Y .
\]
$J$는 대칭행렬이다. 또 각 행의 합이 $y_i-y_i\sum_jy_j=0$인데, 이는 $\sum_jy_j=1$이 항상 성립하기 때문이다.
\end{sol}

% ---------------------------------------------------------------- 6
\item ReLU 계층의 역전파를 다음과 같이 구현한 코드가 있다.
\begin{lstlisting}[numbers=none]
def backward(self, dZ):
    return np.maximum(self.A, 0) * dZ
\end{lstlisting}
무엇이 잘못되었는지 설명하고 올바르게 고쳐라.

\begin{sol}
ReLU $z=\max(0,a)$의 미분은 $a>0$이면 $1$, $a\le0$이면 $0$이다. 따라서 역전파는
upstream gradient $\pd{L}{Z}$를 $a>0$인 곳에서는 \textbf{그대로} 통과시키고 나머지에서는 $0$으로 막아야 한다.

이 코드는 미분값 $1$ 대신 \textbf{순전파의 출력값} $\max(0,a)$를 곱한다.
그래서 $a>0$인 곳의 gradient가 $a$배가 된다. 예를 들어 $a=3$이면 gradient가 $3$배로 부풀고,
$a=0.01$이면 $\frac{1}{100}$로 줄어든다.
연쇄법칙의 국소 미분을 함수값과 혼동한 것이다.
(실습 코드의 \texttt{Relu} 계층은 순전파에서 \texttt{self.A}를 저장하지 않으므로, 이 코드는 그대로는 실행도 되지 않는다.)

순전파에서 저장한 \texttt{mask} $=(A>0)$를 쓰면 된다.
\begin{lstlisting}[numbers=none]
def backward(self, dZ):
    return np.where(self.mask, dZ, 0.0)      # or: dZ * self.mask
\end{lstlisting}
\end{sol}

% ---------------------------------------------------------------- 7
\item 미니배치 학습에서 softmax--CE 계층이 \texttt{(Y - T) / h}를 돌려주는데,
affine 계층의 \texttt{backward}에서도 \texttt{self.dW = self.X.T @ dA / h}처럼 다시 $h$로 나누었다고 하자.
$\pd{L}{W^{(2)}}$와 $\pd{L}{W^{(1)}}$는 정확한 값의 몇 배가 되는가?

\begin{sol}
미니배치 손실 $L=\frac1h\sum_{k}l_k$의 $\frac1h$은 softmax--CE 계층의 \texttt{(Y - T) / h}에서 이미 한 번 들어갔다.
그 뒤의 계층은 연쇄법칙대로 곱하기만 하면 된다.
\begin{itemize}[leftmargin=1.5em,itemsep=2pt]
\item 둘째 affine 계층: \texttt{dA}는 올바른 값인데 \texttt{dW}를 다시 $h$로 나누므로
$\pd{L}{W^{(2)}}$는 정확한 값의 $\boxed{1/h}$배이다.
\item 이 계층이 앞으로 돌려주는 \texttt{dA @ self.W.T}는 나누지 않았으므로 올바르다.
ReLU를 지나 첫째 affine 계층에 도착하는 \texttt{dA}도 올바르다.
그런데 첫째 affine 계층도 같은 클래스이므로 \texttt{dW}를 다시 $h$로 나눈다.
따라서 $\pd{L}{W^{(1)}}$도 정확한 값의 $\boxed{1/h}$배이다.
\end{itemize}
$h=100$이면 두 가중치의 gradient가 모두 $\frac1{100}$이 된다.
$\pd{L}{B}$는 나누지 않았으므로 올바르다.
결과적으로 가중치에 대해서만 학습률을 $\frac{1}{100}$로 줄인 것과 같아서 학습이 매우 느려진다.
gradient 확인(gradient check)을 하면 수치 미분과 정확히 $h$배 차이가 나서 바로 발견할 수 있다.
\end{sol}

% ---------------------------------------------------------------- 8
\item (도전) 손실함수를 제곱 손실 $l=\frac12\sum_i(y_i-t_i)^2$으로 바꾸면
\[
\pd{l}{u_i}=y_i\Bigl((y_i-t_i)-\sum_{j=1}^{n}y_j(y_j-t_j)\Bigr)
\]
임을 보이라. 또 $R=Y-T$라 할 때 이 식을 미니배치 행렬로 쓰면
\begin{lstlisting}[numbers=none]
Y * (R - np.sum(Y * R, axis=1, keepdims=True))
\end{lstlisting}
가 됨을 확인하라.

\begin{sol}
$l$은 $y_1,\ldots,y_n$을 거쳐 $u_i$에 의존하므로 연쇄법칙으로
\[
\pd{l}{u_i}=\sum_{j=1}^{n}\pd{l}{y_j}\pd{y_j}{u_i},\qquad
\pd{l}{y_j}=y_j-t_j .
\]
문제 5에서 $\pd{y_j}{u_i}=\delta_{ij}y_j-y_iy_j$이므로
\begin{align*}
\pd{l}{u_i}&=\sum_{j}(y_j-t_j)(\delta_{ij}y_j-y_iy_j)
=y_i(y_i-t_i)-y_i\sum_{j}y_j(y_j-t_j)\\
&=y_i\Bigl((y_i-t_i)-\sum_{j}y_j(y_j-t_j)\Bigr).
\end{align*}

\textbf{미니배치 행렬.} $Y$, $T$, $R=Y-T$는 모두 $h\times n$이고, $k$번째 행이 $k$번째 데이터이다.
\begin{itemize}[leftmargin=1.5em,itemsep=2pt]
\item \texttt{Y * R}의 $(k,j)$ 성분은 $y_{kj}(y_{kj}-t_{kj})$이다 (성분별 곱).
\item \texttt{np.sum(Y * R, axis=1, keepdims=True)}는 각 행의 합 $\sum_jy_{kj}(y_{kj}-t_{kj})$이고 shape은 $h\times1$이다.
\item \texttt{R - (...)}에서 $h\times1$이 각 행의 $n$개 성분에 broadcast되어 $(k,i)$ 성분이 $(y_{ki}-t_{ki})-\sum_jy_{kj}(y_{kj}-t_{kj})$가 된다.
\item 마지막으로 \texttt{Y *}를 성분별로 곱하면 $(k,i)$ 성분이 위에서 구한 $\pd{l_k}{u_{ki}}$이다.
\end{itemize}
\texttt{keepdims=True}가 없으면 shape이 $(h,)$가 되어 broadcast가 열 방향으로 일어나므로 틀린 값이 된다 ($h\ne n$이면 오류).
미니배치 평균 손실의 gradient는 여기에 $\frac1h$을 곱한 것이며, 실습 코드 \texttt{SoftmaxWithSquaredError.backward}는 마지막에 \texttt{/ h}를 한다.

cross entropy의 $Y-T$와 비교하면, 제곱 손실의 gradient에는 $y_i$가 곱해져 있다.
그래서 정답 class의 확률 $y_i$가 $0$에 가까운 (크게 틀린) 경우에 gradient가 작아져 학습이 느려진다.
\end{sol}

\end{problems}

\end{document}
